\section*{KI-Verzeichnis}

\begin{table}[h]
    \centering
    \begin{tabular}{|>{\raggedright\arraybackslash}p{0.2\textwidth}|>{\raggedright\arraybackslash}p{0.7\textwidth}|}
    \hline
        \textbf{Werkzeug} & \textbf{Beschreibung der Nutzung} \\ \hline
        ChatGPT & Der Einsatz erfolgte in der gesamten Arbeit zur sprachlichen Gestaltung. Konkret diente es dazu, 
        vorhandene Texte stilistisch an den wissenschaftlichen Schreibstil anzupassen und präzise 
        Formulierungen zu entwickeln. Außerdem wurde wie in Kapitel~\ref{methodik1} bzw. Kapitel~\ref{methodik3} 
        beschrieben, die Funktion \textit{ChatGPT Deep Research} für das Identifizieren relevanter Quellen genutzt. 
        Limitationen wurden in besagten Kapiteln reflektiert. Darüber hinaus kam ChatGPT zur Klärung fachlicher 
        Verständnisfragen im Verlauf der Recherche und Umsetzung des Projektes zum Einsatz.\\ \hline
        
        Perplexity & Der Einsatz erfolgte wie in den Kapiteln~\ref{methodik1} und \ref{methodik3} beschrieben, um 
        relevante Quellen zu identifizieren. Wie bei \textit{ChatGPT Deep Research}, wurde der Einsatz in den Kapiteln 
        kritisch hinterfragt und Limitationen identifiziert. \\ \hline
    \end{tabular}
    \caption{Eingesetzte KI-Werkzeuge und Beschreibungen deren Verwendung}\label{tab:ki}
\end{table}